\section{Quantitative separation by logarithmic-length periods}
\label{sec:quantitative-periods}
The complete periodic annihilator was determined in
Theorem~\ref{thm:joint-periodic}. We now estimate the period increments
uniformly in both the radius and the number of collisions. This gives an
explicit separation at large roof frequencies, using three actual
periodic words. Throughout this section $g=\sqrt3-2R$, $d=1/2-R$,
and $e$, $\ell$ and $u$ are the functions in+\eqref{eq:excursion-lengths}. The coordinates below are heights, not
arc coordinates.

Define the half action on $[-d,0]^m$ by
\begin{equation}\label{eq:half-action-v4}
 W_{m,R}(y)=e_R(y_1)-g/2+
       \sum_{j=1}^{m-1}\bigl(\ell_R(y_j,y_{j+1})-g\bigr)+u_R(y_m).
\end{equation}
Reflection symmetry and uniqueness of the full minimizer imply
$E_m(R)=2\min W_{m,R}$. Denote its minimizing coordinates by
$y^{(m,R)}_1,\ldots,y^{(m,R)}_m$.

\begin{lemma}[Uniform Hessian and endpoint estimates]
\label{lem:uniform-half-hessian}
On the entire height box, the Hessian $H=\nabla^2 W_{m,R}$ is
tridiagonal and satisfies
\begin{equation}\label{eq:half-hessian-bounds}
 3I_m\le H\le12I_m,\qquad H_{jj}\le10,\qquad
 \frac12\le-H_{j,j+1}\le2\quad(1\le j<m).
\end{equation}
The last condition is empty when $m=1$. For all $1\le j\le m$,
\begin{equation}\label{eq:height-two-sided}
 \frac3{440}\,20^{-(j-1)}
 \le -y^{(m,R)}_j
 \le \frac1{21}\left(\frac7{10}\right)^{j-1}.
\end{equation}
All constants are independent of $m$ and $R\in I$.
\end{lemma}
\begin{proof}
Write $x(y)=\sqrt{R^2-y^2}$, $a=D/2-x(y)$,
$h=D-x(y)-x(z)$, $v=z-y$, with $D=\sqrt3$.
The elementary bounds on the whole box are
\[
 \begin{gathered}
 x>11/25,\quad |x'|\le5/44,\quad79/200<a<43/100,\\
 79/100<h<86/100,\quad e<11/25,\quad\ell<9/10.
 \end{gathered}
\]
Also $u''=R^2/x^3\ge1/R\ge100/47$ and
$u''\le55225/21296<3$. For the norm of a vector-valued function,
the Hessian is its transverse positive semidefinite quadratic form
plus the Hessian of the vector dotted with its unit direction.
Consequently
\[
 e''\ge(a/e)u''>\frac{79}{88}\frac{100}{47}>\frac32,
 \qquad
 \nabla^2\ell\ge(h/\ell)\operatorname{diag}(u''(y),u''(z))
 >\frac32 I_2.
\]
Each variable of $W_m$ occurs in two such terms, with $u''>3/2$
at the terminal endpoint. This proves the lower bound $H\ge3I_m$,
including $m=1$.

The same norm formula gives
\[
 \ell_{yy},\ell_{zz}
 \le\frac{1+(5/44)^2}{79/100}+\frac{55225}{21296}<4,
 \qquad
 e''\le\frac{1+(5/44)^2}{79/200}+\frac{55225}{21296}<6.
\]
An explicit mixed derivative is
\begin{equation}\label{eq:mixed-length-hessian}
 -\ell_{yz}
 =\frac{(h+vy/x(y))(h-vz/x(z))}{\ell^3}.
\end{equation}
Each numerator factor lies between $78/100$ and $87/100$.
Using $79/100<\ell<9/10$ proves $1/2<-\ell_{yz}<2$.
The diagonal bounds for $W_m$ are therefore $10$ at its first
endpoint, $8$ in its interior, and $7$ at its last endpoint; when
$m=1$ the bound is $9$. Absolute row sums are at most $12$.
The symmetric matrix bound $H\le12I_m$ follows, proving
\eqref{eq:half-hessian-bounds}.

At zero the gradient is $a_0 e_1$, where
$a_0=d/e_R(0)$ satisfies $3/44\le a_0\le1/7$.
Integrate the Hessian on the segment from zero to the minimizer and
write the resulting matrix as $A$. Stationarity gives
$Ay^{(m,R)}=-a_0e_1$. The matrix $B=10I_m-A$ is entrywise
nonnegative and tridiagonal. Its eigenvalues lie in $[-2,7]$, so
$\|B/10\|_2\le7/10$. Hence
\[
 A^{-1}=\frac1{10}\sum_{k\ge0}(B/10)^k.
\]
For the $(j,1)$ entry every power before $j-1$ vanishes. The
nearest-neighbor product at power $j-1$ is at least $20^{-(j-1)}$.
Entrywise nonnegativity and the norm bound yield
\[
 \frac1{10}20^{-(j-1)}
 \le (A^{-1})_{j1}
 \le\frac1{10}\sum_{k\ge j-1}(7/10)^k
 =\frac13(7/10)^{j-1}.
\]
Multiplication by the two bounds on $a_0$ proves
\eqref{eq:height-two-sided}. The matrix is an averaged Hessian of
the actual nonlinear action; no quadratic approximation of the
minimizing orbit has been made.
\end{proof}

\begin{theorem}[Two-sided excess increments]
\label{thm:quantitative-excess}
For every $R\in I$ and $m\ge1$,
\begin{equation}\label{eq:quantitative-excess}
 \frac1{10000}\,400^{-m}
 \le E_{m+1}(R)-E_m(R)
 \le\frac1{300}\left(\frac{49}{100}\right)^{m-1}.
\end{equation}
In particular $E_m$ converges uniformly on $I$ to a continuous
function $E_\infty$, with
\begin{equation}\label{eq:uniform-excess-limit}
 0<E_\infty(R)-E_m(R)
 \le\frac1{153}\left(\frac{49}{100}\right)^{m-1}.
\end{equation}
\end{theorem}
\begin{proof}
For a minimizer $y$ of $W_m$, append the coordinate zero. With
$K(y,z)=\ell(y,z)-g$, the resulting change in half action is
\[
 K(y_m,0)-u(y_m)
 =\sqrt{(g+u(y_m))^2+y_m^2}-(g+u(y_m))
 \le\frac{y_m^2}{2g}.
\]
It follows that $E_{m+1}-E_m\le y_m^2/g$. The upper bound in
\eqref{eq:height-two-sided} and $g>79/100$ give the upper
coefficient $100/34839<1/300$ in \eqref{eq:quantitative-excess}.

Conversely, truncate a minimizer $z$ of $W_{m+1}$ after its first
$m$ coordinates. The removed contribution is
\[
 K(z_m,z_{m+1})+u(z_{m+1})-u(z_m)\ge2u(z_{m+1}),
\]
by \eqref{eq:kernel-coercivity}. Thus
\[
 E_{m+1}-E_m\ge4u(z_{m+1})\ge\frac{2z_{m+1}^2}{R}
 \ge\frac9{45496}\,400^{-m}>\frac1{10000}\,400^{-m}.
\]
Here $u(z)=z^2/(R+x(z))\ge z^2/(2R)$ and
\eqref{eq:height-two-sided} were used. Both comparisons concern
minima of the exact actions. Summing the geometric upper bound
proves \eqref{eq:uniform-excess-limit}. Each $E_m$ is continuous
by Lemma~\ref{lem:excursion-family}; the uniform limit is continuous.
\end{proof}

Let $P_0$ denote the long normal orbit and $P_m$ the $m$th excursion.
For a periodic word $P$ write
\[
 Z_P(u,s,\beta,c;R)
 =\exp i\bigl(u\cdot K_P+sN_P+\beta T_P-ch_P\bigr).
\]
Put $a=\log(100/49)$ and, for $b\ge1$, define
\begin{equation}\label{eq:frequency-period-budget}
 m(b)=1+\left\lceil\frac{\log b}{a}\right\rceil,
 \qquad \gamma=\frac{\log400}{a}-1.
\end{equation}

\begin{theorem}[Explicit large-roof phase separation]
\label{thm:log-period-separation}
For $|\beta|=b\ge1$, all $R\in I$, and arbitrary torus
frequencies $u,s,c$, the three true periods $P_0,P_{m(b)},P_{m(b)+1}$
satisfy
\begin{equation}\label{eq:log-period-separation}
 \max_P|Z_P-1|
 \ge\frac{b\,400^{-m(b)}}{15000\pi}
 \ge\frac{b^{-\gamma}}{2400000000\pi}.
\end{equation}
The exponent is $\gamma=\log400/\log(100/49)-1$.
For a whole band $1\le|\beta|\le B$, the fixed three periods
$P_0,P_{m(B)},P_{m(B)+1}$ suffice, with lower bound
$400^{-m(B)}/(15000\pi)$. Their longest physical word has
$2m(B)+3$ collisions. All assertions are uniform in the radius.
\end{theorem}
\begin{proof}
The actual counts and lengths give the exact identity
\begin{equation}\label{eq:three-period-cancellation}
 Z_{P_{m+1}}\overline{Z_{P_m}}\overline{Z_{P_0}}
       =\exp\bigl(i\beta(E_{m+1}-E_m)\bigr).
\end{equation}
The winding, physical count and induced period cancel. With $m=m(b)$,
\eqref{eq:quantitative-excess} implies
\[
 \frac{b\,400^{-m}}{10000}
 \le b(E_{m+1}-E_m)\le\frac1{300}.
\]
For $|v|\le\pi$, $|e^{iv}-1|\ge2|v|/\pi$. The distance
from one of a product of three unit complex numbers is at most
the sum of their individual distances from one. Applying these
facts to \eqref{eq:three-period-cancellation} proves the first
bound in \eqref{eq:log-period-separation}. Since
$m(b)\le2+\log b/a$, the second bound follows. For the band
assertion use $m(B)$; the upper phase bound still holds, while
$b\ge1$ supplies the printed lower bound. The number of collisions
was proved in Lemma~\ref{lem:excursion-family}.
\end{proof}

\begin{corollary}[An approximate-phase lower bound]
\label{cor:approximate-phase-bound}
Let $q:Y_R^*\to\mathbb C$ have modulus one and be continuous in
neighborhoods of all states in the three cycles used above. Suppose
$q$ is measurable elsewhere and let
\[
 D(q)=\mathop{\rm ess\,sup}_{x\in Y_R^*}
 \left|e^{i(u\cdot\kappa_R^*(x)+s r_R^*(x)+\beta\tau_R^*(x))}
                 q(F_R^*x)-e^{ic}q(x)\right|.
\]
For $b=|\beta|\ge1$ and $m=m(b)$,
\begin{equation}\label{eq:approximate-phase-bound}
 D(q)\ge\frac{b\,400^{-m}}{10000\pi(m+1)}.
\end{equation}
\end{corollary}
\begin{proof}
The return branches through these cycles are regular, and their
section measure has full support. Continuity therefore extends the
essential bound to every state of each selected cycle. Multiplication
around a cycle of induced period $h$ telescopes the unit phases and
gives $|Z_P-1|\le hD(q)$. In
\eqref{eq:three-period-cancellation} the sum of the three periods is
$1+m+(m+1)=2(m+1)$. The lower sine bound in the preceding proof
therefore gives
$2b(E_{m+1}-E_m)/\pi\le2(m+1)D(q)$. Apply the lower bound
in \eqref{eq:quantitative-excess}.
\end{proof}

The estimates are quantitative statements about the specified mechanical
return map, and in particular about continuous circle-valued approximate
phases. They do not assert regularity of a merely measurable transfer
function, a lower bound for a spectral eigenfunction on these orbits,
or a transfer-operator resolvent estimate. Those are distinct analytical
steps when using \eqref{eq:approximate-phase-bound} in a high-frequency
argument. The raw characteristic-function tail is not replaced by the
period discrepancy.
